\documentclass[a4paper]{article}
\usepackage{graphicx}
\usepackage{parskip}
\usepackage[margin=1in]{geometry}
\usepackage{url}

\title{Assignment 3 Response Sheet}
\date{2026 Semester 1}

\begin{document}

\maketitle

\section{Instructions}
\textbf{Submission deadline: 26 April 2026 at 23:00. (Week 8, Sunday)}
Please type your answers in the spaces provided. Do not repeat the questions in your answers. All tasks are compulsory.

Note that you must correctly cite and reference any sources you have consulted. You may use any internationally recognized referencing style such as APA or IEEE referencing.

Submit this document to the ``Assignment 3 Main Submission'' inbox by the published deadline.

Type your student ID number in the space provided below.

\textbf{Student ID number: 520465817}

\section{Task 1}
\subsection{Introduction}
Large language model (LLM) agents are increasingly designed as systems that combine reasoning, action, memory, and external capability use rather than as standalone text generators. Early work on retrieval-augmented generation showed that language models can improve performance by accessing external non-parametric memory at inference time rather than relying solely on knowledge stored in model parameters \cite{lewis2021rag, izacard2022atlas}. Later work on tool use extended this idea from document retrieval to action execution by enabling models to call external functions and APIs \cite{schick2023toolformer, patil2023gorilla, qin2023toolllm}. More recently, the literature has begun to treat agent skills as a distinct abstraction: reusable procedural capabilities that may package instructions, scripts, workflows, tools, and task-specific resources into units that can be invoked across tasks \cite{jiang2026procedural, jiang2026sok, anthropic2025skills}. Unlike atomic tools, which typically expose a specific callable function, skills operate as higher-level procedural modules. This shift is important because it changes the retrieval problem. If the objects being selected are not atomic tools but higher-level procedures, then retrieval quality depends not only on semantic relevance but also on procedural suitability.

The research problem addressed in this review is therefore not simply whether large capability libraries are difficult for agents to manage. That broad point is already well established. The more specific issue is whether current methods can distinguish between skills that appear semantically similar at the level of names or descriptions, yet differ in preconditions, workflow steps, dependencies, outputs, or side effects. This distinction matters particularly in personal-agent settings, where users may accumulate many overlapping skills over time without maintaining a carefully curated capability library. In such settings, flat skill descriptions may create ambiguity even when the correct skill exists in the library. A retrieval system may therefore fail not because a capability is missing, but because the representation and selection process do not preserve procedural differences strongly enough.

This literature review is organised around that problem. Rather than summarising papers one by one, it examines the literature through the questions most relevant to scalable skill retrieval. First, why should skills be treated as a distinct retrieval problem rather than as a simple extension of tool selection? Second, what does the emerging skill literature reveal about the scaling problem itself? Third, what representations and retrieval mechanisms have been proposed for skills, and what limitations remain? Finally, what evaluation practices are currently used, and do they adequately test semantic confusability, context efficiency, and downstream task performance? The review therefore compares prior work not only in terms of whether it improves overall agent behaviour, but also in terms of what kind of skill representation it assumes, how retrieval is performed, and how convincingly it is evaluated under scale. The review argues that the literature already provides strong evidence that flat or weakly structured skill selection becomes unreliable as skill libraries grow, and that structure-aware approaches are promising. However, it also argues that existing work does not yet provide a clear cross-family comparison of which representational properties most effectively distinguish semantically similar but procedurally different skills while remaining retrieval-efficient. This gap motivates the research questions stated at the end of the review.

\subsection{Conceptual Dimensions of Skill Retrieval}
A useful conceptual distinction is needed before comparing the literature. A \emph{skill artifact} is the actual stored capability, such as a prompt module, executable program, or folder of instructions and scripts. A \emph{selection representation} is the form in which that skill is made available for retrieval or routing, such as a short textual description, a structured schema, a hierarchy position, or a graph node. A \emph{retrieval policy} is the mechanism that selects candidates, such as direct prompt exposure, dense retrieval, hierarchical traversal, graph exploration, or hybrid reranking. Finally, an \emph{evaluation protocol} specifies what success means: retrieval accuracy, downstream task success, context efficiency, robustness under scale, or some combination of these dimensions. This distinction matters because many systems are discussed as if artifact, representation, and retrieval policy were the same thing. In practice, however, a skill may be stored as code or documentation while being represented for selection through only a short textual descriptor. Similarly, a system may use a hierarchical index without changing the internal structure of each skill. Making these layers explicit allows a clearer comparison of what prior work actually contributes.

\subsection{Why Skills Are a Distinct Retrieval Problem}
The first important question is whether skills are genuinely different from tools for the purposes of retrieval. The recent skill literature strongly suggests that they are. The survey \emph{Agent Skills from the Perspective of Procedural Memory} conceptualizes skills as reusable procedural memories that enable an agent to reuse problem-solving procedures rather than repeatedly reasoning from scratch \cite{jiang2026procedural}. This framing is useful because it highlights the procedural nature of skills. The survey organizes the field into acquisition, representation, invocation, and refinement, showing that skills can be encoded as prompts, programs, APIs, or more complex procedural modules. Its main strength is to clarify that a skill is not just a callable function, but a reusable unit of procedure. However, the paper is primarily taxonomical. It does not provide a benchmark or retrieval method that shows how agents should select among many semantically overlapping skills. Its importance for the present thesis lies in problem framing rather than direct methodological evidence.

A similar point is developed more strongly in the \emph{SoK: Agentic Skills --- Beyond Tool Use in LLM Agents} \cite{jiang2026sok}. That paper argues that skills extend beyond tools, plans, or memory because they operate as executable procedural modules that can encapsulate complex task knowledge. This is directly relevant to the present topic because it implies that retrieval methods designed for atomic tools may not transfer cleanly to skills. If tools are primarily selected on the basis of function-level relevance, while skills must be selected on the basis of procedural appropriateness, then the retrieval target itself changes. The value of the SoK paper is therefore conceptual: it strengthens the argument that skill retrieval should not be treated as merely a larger version of tool retrieval. However, like the procedural-memory survey, it does not yet solve the practical problem of large-scale skill selection.

Early skill systems further show why this distinction matters. SkillAct introduces textual skill abstractions represented as a tuple consisting of a name, natural-language instructions, and example trajectories, all appended directly to the prompt of a baseline agent \cite{liu2024skillact}. The paper demonstrates that skill abstractions can improve performance, especially by reducing execution errors in ALFWorld. Its contribution is important because it provides an early and transparent baseline for skill-augmented agents. However, its retrieval model is effectively absent: the available skills are exposed as a flat prompt list. This means the system works only because the experimental setting is small enough for full prompt inclusion to remain manageable. The paper therefore supports the claim that skills are useful, but also exposes the weakness of flat prompt-level skill access when libraries become large or internally redundant.

Voyager provides a stronger architectural example of a skill library \cite{wang2023voyager}. It stores successful behaviors as executable code and indexes them by embeddings of their natural-language descriptions. This is significant because it shows that skills can function as reusable external capabilities rather than one-off generated plans. The paper demonstrates that a growing skill library improves exploration and transfer in Minecraft. However, its retrieval mechanism remains shallow: skills are recalled through embedding similarity over descriptions rather than through procedural structure. As a result, Voyager is best understood as a strong early precursor to skill-library architectures rather than a complete answer to scalable skill retrieval. It demonstrates the utility of externalized skills, but does not study what happens when multiple skills share similar language while differing in operational logic.

\subsection{Evidence That Skill Scaling Is a Real Problem}
The next major theme in the literature concerns whether the scaling problem is real and, if so, what exactly causes it. Three recent papers are especially important here: \emph{SkillsBench}, \emph{When Single-Agent with Skills Replace Multi-Agent Systems and When They Fail}, and \emph{Agent Skills: A Data-Driven Analysis of Claude Skills} \cite{li2026skillsbench, li2026singleagent, ling2026claudeskills}. Together, these papers shift the discussion away from generic claims about capability growth and toward more specific evidence about skill-library behavior under scale.

SkillsBench evaluates whether skills actually help LLM agents across diverse tasks by comparing conditions with no skills, curated skills, and self-generated skills \cite{li2026skillsbench}. Its most useful finding is not merely that curated skills improve performance, but that smaller and more focused skill sets often outperform larger and more comprehensive ones. This matters because it suggests that the effectiveness of a skill library depends not only on the presence of reusable procedures, but also on how tightly they are scoped and represented. The benchmark is valuable because it tests skill usefulness across domains rather than assuming that more skills are always better. However, it only weakly addresses the retrieval problem. The relevant curated skills are effectively assumed to be available, so the study does not directly evaluate how the correct skill should be retrieved from a large ambiguous library. Its contribution is therefore diagnostic rather than decisive: it shows that skill packaging matters, but not how scalable skill selection should be performed.

The paper \emph{When Single-Agent with Skills Replace Multi-Agent Systems and When They Fail} is even more directly relevant \cite{li2026singleagent}. It formalizes skills as descriptor-policy-backend tuples and studies how skill selection behaves as the number of available skills increases. Its most important contribution is the finding that performance remains relatively stable until a threshold is reached, after which selection accuracy drops sharply. Crucially, the paper shows that this degradation is driven not merely by the number of skills, but by semantic confusability among them. This is one of the strongest direct pieces of evidence for the present thesis. It isolates semantic overlap as the core failure mode of flat skill selection. At the same time, the paper is weaker as a solution paper. Its proposed hierarchical mitigation is intentionally simple and does not determine which hierarchy design or retrieval structure is best. Nor does it represent skills through rich procedural structure. The paper therefore motivates the research problem very strongly, but does not yet solve it.

The data-driven analysis of Claude skills contributes a different kind of evidence \cite{ling2026claudeskills}. Rather than proposing a new retrieval method, it measures the properties of a large public skill ecosystem, including redundancy, length variation, and supply-demand imbalance. This is highly relevant because it shows that real skill libraries are already noisy, redundant, and unevenly structured. The paper is particularly useful for context-efficiency arguments, since it shows that while many skills are short, a non-trivial minority are long enough to create prompt-budget pressure if loaded indiscriminately. It also highlights that discovery often happens through lightweight metadata before full instructions are loaded. However, the paper is not a retrieval-comparison study. Its contribution is ecosystem diagnosis rather than direct evidence about which retrieval methods work best. Still, it supports the claim that scalable discovery and selective loading are practical necessities rather than theoretical luxuries.

\subsection{Flat and Weakly Structured Skill Representations}
Having established that skills are distinct from tools and that large skill libraries create real confusability problems, the next question is how existing work represents skills for selection. The simplest family of representations remains flat or weakly structured textual skills, as seen in SkillAct and parts of Voyager \cite{liu2024skillact, wang2023voyager}. These approaches work because they make reusable procedures accessible to the model with minimal infrastructure. Their strengths are simplicity, interpretability, and ease of authoring. Their weakness is that the differentiating information remains embedded in free-form text. A flat description can indicate what a skill is roughly for, but it may not clearly encode conditions under which one apparently similar skill should be preferred over another. This limitation becomes especially serious when the representation available to the retriever is shorter and shallower than the full skill artifact. Flat prompt-based skill systems therefore make a valuable baseline, but not a scalable endpoint.

\subsection{Structure-Aware Skill Representations}
A more promising direction is to introduce structure explicitly into the skill representation itself. One example is \emph{Learning Agent Skills from Demonstrations and Generating Knowledge Graph}, which decomposes complex demonstrations into meta-skills, generates behavior trees, and transforms the resulting information into a knowledge graph \cite{learningkg2026}. This paper is relevant because it shows that skills can be represented through explicit internal structure rather than only through free-form descriptions. By moving from trajectories to behavior trees and graph representations, it makes the procedural nature of skills more visible and interpretable. However, the paper is heavily tailored to robotics assembly domains and relies on a predefined meta-skill vocabulary. Moreover, the graph serves more as an organizational and visualization layer than as a thoroughly evaluated large-scale retrieval mechanism. Its main value for the present review is therefore to support the general idea that skills can be decomposed into richer procedural structures.

\subsection{Hierarchical Skill Retrieval}
Two more directly relevant large-scale skill infrastructure papers are AgentSkillOS and SkillNet \cite{li2026agentskillos, liang2026skillnet}. AgentSkillOS organizes skills into a hierarchical capability tree and then uses this structure for retrieval and DAG-based orchestration. Its strongest contribution is empirical: it demonstrates that structured retrieval can approximate oracle skill selection and that structured orchestration improves downstream task performance compared with flat invocation. This makes it one of the most important system-level papers for the present thesis. However, the capability tree assumption is also its main limitation. Real skills often belong to multiple overlapping functional categories, and a strict tree risks brittle routing decisions. Moreover, the evaluation focuses on overall system performance and orchestration quality rather than isolating retrieval accuracy under highly confusable skill competitors. AgentSkillOS therefore provides strong evidence that structure helps, but weaker evidence about which structural properties are specifically responsible for disambiguating semantically similar skills.

\subsection{Graph-Based and Relational Skill Organization}
SkillNet takes a broader relational approach \cite{liang2026skillnet}. It introduces a taxonomy layer, a relation graph with typed edges such as similarity, composition, and dependency, and a package layer for deployment. This paper is particularly relevant because it refuses to treat skills as isolated items and instead models them as connected assets within a structured network. This is conceptually closer to the needs of large skill ecosystems, where similarity and compositionality often coexist. The main limitation is that the paper functions more as a full infrastructure proposal than as a clean retrieval-comparison study. It provides evidence that structured skill libraries improve downstream performance, but offers less fine-grained analysis of retrieval accuracy, context cost, or robustness under semantic confusability. It therefore strengthens the case for relational structure, but stops short of determining how such structure should be exploited for runtime disambiguation.

Hierarchical retrieval deserves separate attention because it is the clearest structure-aware response to scaling so far. Both \emph{When Single-Agent with Skills Replace Multi-Agent Systems and When They Fail} and AgentSkillOS suggest that coarse-to-fine routing can mitigate the collapse of flat selection under scale \cite{li2026singleagent, li2026agentskillos}. The core intuition is straightforward: forcing the agent to choose directly from a large flat pool is harder than guiding it through a sequence of smaller local decisions. This can reduce both the search space and the prompt burden. However, hierarchy introduces its own assumptions. It assumes that an early routing decision can be made reliably and that the correct skill is reachable through a small number of functional branches. When skills overlap across multiple categories, or when the hierarchy itself is derived from semantics rather than procedure, these assumptions can fail. Thus, hierarchical skill retrieval appears promising as a scalability aid, but it is not yet clear whether it resolves the deeper procedural disambiguation problem or merely postpones it.

Graph-based and relational organization provide a different response to that problem. By allowing multi-parent and typed relationships, graphs potentially model the reality of skill ecosystems more faithfully than strict trees. SkillNet is the clearest direct skill example, while the knowledge-graph paper on learned skills from demonstrations shows that structured relational representations are at least plausible at the skill level \cite{liang2026skillnet, learningkg2026}. The main promise of graph-based organization is that it can encode not only similarity but also composition, dependence, and workflow compatibility. However, graph approaches also face a major difficulty: graph quality depends on how nodes and edges are constructed. If similarity edges are derived from the same shallow semantics that already confuse flat retrieval, then the graph may reproduce the original ambiguity rather than solve it. This means that graph structure is not automatically beneficial; its utility depends on whether the encoded relations capture meaningful procedural distinctions.

\subsection{Transferable Ideas from Tool and Document Retrieval}
At this point it is reasonable to ask whether ideas from the tool-selection and document-retrieval literature can help. They can, but only as \emph{transferable methodological evidence}, not as direct proof about skill retrieval. For example, SciToolAgent shows that graph-based organization of tools can improve planning and multi-step execution by representing input-output formats, dependencies, and safety relations explicitly \cite{ding2025scitoolagent}. This paper is useful because it demonstrates the value of dependency-aware structured retrieval. However, its unit of organization is the tool, not the skill, and its graph is manually curated in a highly specialized domain. Similarly, LATTICE demonstrates that hierarchical traversal with calibrated path scoring can outperform flat retrieve-then-rerank pipelines for document retrieval \cite{gupta2025lattice}. This is methodologically interesting because it shows that LLMs can be used to navigate an index rather than only rerank retrieved candidates. Yet its direct transfer to skills is limited because documents are not procedural capability units. ToolLLM and Gorilla further establish the practical importance of retrieval when capability spaces become large, but they remain fundamentally tool-centric and do not address procedural skill structure \cite{qin2023toolllm, patil2023gorilla}. These papers should therefore be used cautiously: they provide ideas for hierarchy, graph structure, reranking, and representation design, but not direct evidence that those methods solve skill retrieval.

\subsection{Evaluation Practices and Unresolved Gaps}
One of the clearest gaps in the literature emerges when evaluation practices are examined closely. Many papers report overall task success, but fewer isolate retrieval quality itself. SkillsBench measures whether skills help, but not large-library disambiguation \cite{li2026skillsbench}. AgentSkillOS and SkillNet demonstrate strong system-level gains, but retrieval, orchestration, and planning improvements are partly entangled \cite{li2026agentskillos, liang2026skillnet}. The Claude-skills analysis reveals redundancy and selective-loading pressure, but does not compare retrieval methods experimentally \cite{ling2026claudeskills}. The Li paper on single-agent skill failure directly targets semantic confusability, which makes it especially important, yet its mitigation strategy is intentionally simple \cite{li2026singleagent}. Across these works, three recurring evaluation limitations appear. First, semantically confusable but procedurally distinct skills are not always tested explicitly. Second, context efficiency is often discussed qualitatively rather than measured directly in tokens, shortlist size, or cost. Third, some structured approaches are compared against relatively weak baselines, making it difficult to isolate how much value comes from structure-aware selection itself rather than broader system engineering improvements.

These limitations matter because the present research problem is more specific than generic retrieval quality. A method may retrieve useful skills on average and still fail on the exact cases that matter most for a personal skill library: cases where surface similarity is high but the correct procedure is different. Likewise, a method may improve overall task performance through better orchestration without demonstrating that the representation itself improves retrieval under confusability. The strongest literature therefore points in a common direction without fully closing the question. Skills are useful. Large skill libraries create real ambiguity. Structure-aware organization appears promising. Yet the literature still does not clearly determine which representation best preserves procedural distinguishability while remaining efficient and effective at runtime.

\subsection{Synthesis and Research Questions}
The literature therefore suggests that flat skill selection becomes unreliable as skill libraries grow and semantic confusability increases. Although structure-aware approaches are promising, it remains unclear which representational properties most effectively preserve procedural distinguishability while remaining context-efficient and useful in downstream execution. This gap motivates the following research questions:

\textbf{RQ1:} How can agent skills be represented and retrieved so that agents can accurately select among semantically similar but procedurally distinct skills as skill libraries scale?

\textbf{RQ2:} To what extent do structure-aware skill representations improve retrieval accuracy, context efficiency, and downstream task performance compared with flat description-based skill retrieval?

These questions arise naturally from the literature rather than being imposed on it. Existing work has already shown that flat selection becomes unreliable under scale, that the problem is linked to semantic confusability, and that hierarchical or graph-based organization may help. What remains unresolved is which representational properties most effectively preserve procedural differences between competing skills, and how that effect should be measured. The remainder of the project is therefore best understood not as a generic study of agent scaling, but as a focused investigation into structure-aware skill retrieval under semantic confusability.

\section{Task 2}
The research problem addressed in this project is that current agent skill selection approaches often rely on flat natural-language descriptions or general semantic similarity, which may be insufficient for distinguishing semantically similar but procedurally distinct skills. As the literature review shows, this is not simply a problem of large libraries becoming too long for the prompt. Rather, it is a representational and retrieval problem. Skills are higher-level procedural units that can encode workflows, dependencies, preconditions, tool usage, and execution constraints. When such units are represented only through shallow text descriptions, the retrieval mechanism may rank the wrong skill even if the correct one is present. This can reduce retrieval precision, increase prompt overhead through unnecessary candidate loading, and ultimately harm downstream task performance. Although existing work has explored hierarchical and graph-based organization, it remains unclear which forms of structure most effectively preserve procedural distinguishability in large skill libraries.

This project aims to make three contributions. First, it will clarify the distinction between \emph{skill artifacts}, \emph{selection representations}, and \emph{retrieval policies}. This conceptual separation is useful because many existing discussions conflate the skill itself with the representation exposed for retrieval, or conflate representational structure with the retrieval method used to navigate it. Second, the project will construct or adapt an evaluation setting that stresses semantic confusability between procedurally different skills. Rather than evaluating retrieval only on broad relevance, it will target the exact failure mode identified in the literature: skills that look similar at the description level but require different procedures in practice. Third, the project will empirically compare flat and structure-aware skill representations on retrieval accuracy, context efficiency, and downstream task performance in order to identify which structural properties most improve scalable skill selection.

The proposed methodology is comparative rather than purely constructive. The project will evaluate multiple representation-and-retrieval conditions arranged as a small experimental ladder. The first baseline will be \emph{full in-context exposure}, in which all available skill metadata is presented directly without external retrieval. This condition provides a simple reference point and captures the flat prompt-selection regime used by early skill systems. The second baseline will be \emph{flat description-based retrieval}, in which each skill is represented primarily as a text item and retrieved through semantic similarity. This condition is important because it reflects the most natural non-structured solution and serves as the strongest flat baseline. The structure-aware conditions may include one or more of the following: a schema-based representation that explicitly surfaces fields such as purpose, preconditions, inputs, outputs, and workflow steps; a hierarchical representation that organizes skills into a capability tree or similar coarse-to-fine routing structure; and a graph-based representation that encodes typed relations such as similarity, dependency, and composition. If time constraints limit implementation breadth, the project can still remain credible by comparing the two flat baselines against one strong structure-aware alternative.

A critical part of the project is the benchmark design. The literature suggests that existing evaluations do not sufficiently stress semantically confusable skills, so the benchmark should be aligned directly with that gap. The proposed benchmark will use controlled clusters of skills that have high surface similarity but meaningful procedural divergence. A valid confusable cluster should satisfy three conditions. First, the skills should overlap strongly in name, description, or trigger phrases, so that a flat retriever or LLM could plausibly confuse them. Second, the skills should differ in one or more operationally significant dimensions, such as preconditions, required tools, workflow steps, input and output artifacts, side effects, or verification criteria. Third, wrong selection should matter: choosing the wrong skill should produce a different workflow, incorrect artifact, missing checks, or otherwise degraded task success. Candidate clusters may include distinctions such as code review versus pull-request review versus review-comment resolution, or changelog drafting versus release workflow versus patch-release procedure. The point is not to rely only on raw public skills, but to construct benchmark units that directly test the present research problem.

The evaluation will operate at both retrieval and downstream levels. Retrieval quality will be measured using metrics such as top-1 accuracy, top-k recall, and ranking quality. These metrics are necessary because a method may appear effective overall while still failing to rank the correct skill highly under confusability. Context efficiency will be measured by the amount of skill information exposed during selection, such as prompt length, token cost, or candidate shortlist size. This is important because structure-aware methods may succeed partly by narrowing the search space and reducing the burden of exposing the full skill library. Downstream task performance will be measured by whether the agent successfully completes the task after the selected skill or shortlist is supplied. This final layer is essential because retrieval improvements are only valuable if they translate into better behavior. Additional experiments should vary library size and semantic overlap so that the project can observe how each method behaves as the number of competing skills grows. If feasible, qualitative failure analysis should also be included in order to identify whether errors arise from weak representation, routing mistakes, or downstream execution problems.

This evaluation design is convincing for two main reasons. First, it directly targets the failure mode identified in the literature rather than an easier surrogate problem. Instead of asking only whether retrieval works in general, it asks whether retrieval still works when semantically similar skills must be disambiguated on procedural grounds. Second, it separates dimensions that are often entangled in prior work. By measuring retrieval accuracy, context efficiency, and downstream task success separately, the study can determine whether a structure-aware method improves retrieval itself, merely reduces prompt burden, or produces benefits only after orchestration and execution. In this sense, the evaluation is designed not only to produce comparative results, but also to support a persuasive argument about why those results matter.

Overall, the proposed project is a direct response to the literature review. Task 1 showed that existing work has diagnosed flat-skill selection failure and proposed promising structured alternatives, but has not yet clearly established which representations best preserve procedural distinguishability. Task 2 responds by proposing a controlled empirical comparison of flat and structure-aware skill representations under the exact conditions where current methods appear weakest. If successful, the project will contribute both a clearer conceptual framework for analyzing skill-selection systems and empirical evidence about how scalable personal skill libraries should be organized and retrieved.

\newpage
\bibliographystyle{IEEEtran}
\bibliography{references}

\end{document}
